\section{Why pooling bounds a corrupted record}\label{sec:mechanism}

Consider a record whose covariate was entered wrongly, for example multiplied by four. Its linear
predictor becomes extreme, so its predicted risk goes towards zero or one and it sorts into an extreme
risk group. Its outcome, however, was generated by the covariate the patient actually had. If that
outcome contradicts the recorded prediction, the record is a bad leverage point in the sense of
\citet{rousseeuwleroy1987}; if it agrees, it is a good one. This section shows what one such record can
do to a goodness-of-fit statistic, and why pooling limits it.

\subsection{The record in the fit and in the test}

The record's contribution to the likelihood score is $x_i(y_i-\hat\pi_i)$ and its contribution to the
information is $\hat\pi_i(1-\hat\pi_i)x_ix_i^\top$. At an extreme prediction the second is close to
zero. The first is close to zero when the outcome agrees with the prediction and of the order of
$\lVert x_i\rVert$ when it contradicts it, so a contradicting record pulls the fit, and with one such
record at an arbitrarily extreme covariate the logistic maximum-likelihood estimator breaks down
\citep{croux2002breakdown}. In our designs a covariate multiplied by four changes the fitted slope by
less than $5\%$ with ten corrupted records in $1000$; a sign error, which reverses the prediction and so
makes most corrupted records bad leverage points, changes it by $14\%$.

A corrupted record sits alone at an extreme of the predicted risk. A test is fragile when it divides
that record's contribution by something the record, or the refit it causes, can make small. One such
divisor is the record's own fitted variance $\hat\pi_i(1-\hat\pi_i)$, which vanishes at an extreme
prediction. Another is the information that the test's direction keeps after the fit: a smooth
direction close to the fitted linear predictor keeps little, so a score test that divides by it
magnifies whatever the corrupted record and the slightly moved fit leave along that direction. We call
tests that standardise each record separately \emph{record-level}; Stukel's score test, the cubic
calibration test and the GiViTI belt are of this kind. A test is \emph{pooled} when it combines each record's residual with those
of other records, over groups of risk, over a learned partition or over neighbourhoods in covariate
space, and standardises the sum by the variance of the whole pool.

In the grouped residual \eqref{eq:resid} the corrupted record acts in two ways. Let $\theta$ be the
probability that its outcome contradicts its recorded prediction. In our designs, for a covariate
multiplied by four, $\theta$ is about $0.3$ over all corrupted records and $0.1$ among those that reach
the top group; for a sign error it is about $0.7$ and $0.8$. The record shifts the
numerator: at the top of the risk scale, where $\hat\pi_i\approx1$, its expected contribution to
$O_g-E_g$ is about $-\theta$, and at the bottom about $+\theta$. It also leaves $V_g$ understating the
variation it brings, because $V_g$ is computed at $\hat\pi_i(1-\hat\pi_i)\approx0$ while the record
varies with its true risk. With $m_c$ corrupted records among the $m$ records of a group, $\bar v$ the
mean fitted variance of a clean member and $\bar v_c$ the mean true variance of a corrupted one, a
heuristic first-order approximation gives
\begin{equation}\label{eq:twomoments}
  \mathbb{E}[r_g]\;\approx\;\mp\frac{m_c\,\theta}{\sqrt{(m-m_c)\,\bar v}},
  \qquad
  \operatorname{Var}(r_g)\;\approx\;1+\frac{m_c\,\bar v_c}{(m-m_c)\,\bar v},
\end{equation}
with the minus sign in the top group and the plus sign in the bottom one. In both expressions the
corrupted records are in the numerator and the clean groupmates form the whole denominator: the
contribution of one corrupted record among twenty-four clean ones is divided by the square root of
their summed variance.

\subsection{Bounded influence at fixed coefficients}

\begin{proposition}[Bounded influence in the covariates]\label{prop:bounded}
Hold the coefficients at $\hat\beta$. Replace one record's covariate by an arbitrary value $x^{*}$,
leaving its outcome and every other record unchanged, and write $\eta^{*}=x^{*\top}\hat\beta$. For an
ungrouped directed score test with direction $z(\eta)=\eta|\eta|$ the statistic grows without bound as
$\eta^{*}$ moves in the direction in which the record's outcome contradicts its prediction: the record's contribution to the score is of order $\eta^{*2}$, while its contribution to
the information, $z(\eta^{*})^{2}\hat\pi^{*}(1-\hat\pi^{*})$, tends to zero. For the grouped statistic,
if the record stays in its group and $V_g>1/4$,
\[
  |\Delta r_g|\;\le\;\Bigl(V_g-\tfrac14\Bigr)^{-1/2}
  \Bigl\{1+\frac{|r_g|\sqrt{V_g}}{8\,(V_g-\frac14)}\Bigr\},
\]
and if it changes group the same bound holds with $2$ in place of $1$ for every group it touches.
None of these bounds involves $\eta^{*}$.
\end{proposition}

The proof is in Appendix~\ref{app:proofs}. The proposition is exact in external validation, where
the coefficients are frozen and nothing is refitted. When the model is refitted to the contaminated
data it bounds the record's direct influence on the statistic, not its influence through
$\hat\beta$, which is that of the maximum-likelihood estimator; the measured slope changes above say
how large that is in our designs. This is a finite-sample bounded-influence property in the spirit of
\citet{hampel1986robust}, not a statement about the asymptotic influence function.

The bound is small only when $V_g$ is large. At the default partition an extreme group holds about
twenty-five records whose mean risk is near zero or one, so $V_g$ can be close to $1/4$ and the bound is of
order one; that is where the default partition loses its protection (Section~\ref{sec:severity}), and
ten groups, with $V_g$ in the tens, restore it. The bound holds for any statistic built from grouped
residuals, the Hosmer--Lemeshow statistic included; EDGE inherits it and spends the grouped residuals on
a few directions with unit weights, so that nothing in its denominator depends on the fit's remaining
information. Bounded directions alone do not give the bound its force: Proposition~\ref{prop:bounded}
concerns a record's direct contribution at fixed coefficients, and a score test for the same cubic
directions in the predicted risk passes that test too, yet once the model is refitted it is as fragile
as Stukel's test, and so is EDGE's own score form (Section~\ref{sec:whatprotects}).

Any form of pooling supplies the divisor: le Cessie's test pools over neighbourhoods in covariate
space and is protected (Section~\ref{sec:contamination}). For grouping on the risk, the protection is
set by how many clean records share the groups where corrupted records land, and so by the
partition rather than the contamination rate. The Hosmer--Lemeshow statistic with equal-width rather
than equal-frequency risk groups shows this directly: its extreme groups hold few records, and one
corrupted record in a thousand raises its false-alarm rate to $0.151$ (Table~\ref{tab:severity}). The
protection ends when corrupted records fill a group, because as $m_c$ approaches $m$ the denominator
empties. The heuristic describes the cubic basis well; how the symmetric basis responds to the
partition under a sign error, which Section~\ref{sec:severity} reports, it does not explain.

The bound concerns records whose \emph{prediction} is wrong. A record whose \emph{outcome} is wrong,
for example an event that was never recorded, changes $O_g$ in the direction a miscalibration would,
and a calibration test should detect it.
